\subsection{Related work}\label{sec:relwork}

\paragraph{The expressiveness question.}
For historical background, we refer to the surveys by \citet{pratt1992},
\citet{maddux1991}, and \citet{burris-legris}, the monograph of
\citet{brady2000}, and the historical notes of \citet{tarski-givant1987}.
\CoR\ is older than first-order logic. De Morgan introduced composition
and converse in his papers on the syllogism of 1850 and
1860~\citep{demorgan1856,demorgan1864}, Peirce developed the algebra
through the 1870s and 1880s~\citep{peirce1870,peirce1880}, and Schr\"oder
devoted an entire volume to it in 1895~\citep{schroeder1895}. In contrast,
first-order logic was formalized only in Hilbert's lectures of 1917--18 and
in Hilbert and Ackermann's 1928
book~\citep{hilbert-lectures,hilbert-ackermann1928,moore1988}. In the
terminology of the time, to \emph{condense} a quantified statement about
relations meant eliminating its quantifiers to express it within the
calculus~\citep[p.~233]{loewenheim1915}---the translation of
Proposition~\ref{prop:upper-translation} under its older name. Schr\"oder
claimed in 1895 that every such statement could be
condensed~\citep[p.~551]{schroeder1895}. Korselt found a counterexample,
published by L\"owenheim in 1915~\citep{loewenheim1915}, which lies
exactly on the variable boundary studied here: ``the domain has at most two
elements'' requires three variables (\FOthree) and can be condensed,
whereas ``the domain has at most three elements'' requires four (\FOfour) and
cannot. Statements in \FOfour\ were thus known to lie outside \CoR\
before first-order logic existed as a formal system, while the status of
\FOthree\ remained open.
Tarski's 1941 paper presented further
\FOfour\ statements outside \CoR,
observed that the calculus remained at essentially the same stage of development as forty-five years earlier,
and posed the problem of characterizing the
condensable sentences~\citep[pp.~74, 89]{tarski1941}.
According to Tarski and
Givant, Tarski obtained the \FOthree\ characterization in the early 1940s
and presented it at a 1945 Berkeley seminar. It was first published with full
proofs in their 1987 monograph~\citep[pp.~xvii, 88--89]{tarski-givant1987},
and Maddux provided a modern treatment in his monograph~\citep{maddux2006}.
This settled the expressiveness question and forms the basis for
Proposition~\ref{prop:upper-translation} and Corollary~\ref{cor:gap},
but the computational cost of this translation was not examined.

\paragraph{The succinctness question.}
Expressive equivalence, however, implies nothing about formula size. Tarski
and Givant defined their translation recursively, providing no bound on its
growth~\citep[\S3.9]{tarski-givant1987}. The folklore translation behind
Proposition~\ref{prop:upper-translation} applies bottom-up DNF conversion to
the quantifier-free Boolean shell of a formula---the propositional normal-form
theorem of \citet{hilbert-ackermann1928}, read as a conversion procedure by
\citet{kleene1952} and \citet{schoening1989}---and then applies
\eqref{eq:basic-identity} to each disjunct. This bottom-up conversion is
the sole source of the exponential blow-up. In contrast, the reverse
translation is linear: structural induction turns a \CoR\ term of size $m$
into an \FOthree\ formula of size $O(m)$, so the blow-up occurs in one
direction only. Whether this exponential blow-up is unavoidable has remained
open since the equivalence was established. The query $\Psi_k$ is due to
\citet{nakamura2022}, who showed that it has no subexponential translation
into \emph{positive} \CoR\ (the fragment without complement), leaving the
case of full \CoR\ open. Theorems~\ref{main:omega} and~\ref{main:phi} settle
this question.

\paragraph{Succinctness of logics and lower-bound methods.}
Succinctness is a standard measure for comparing logics of the same
expressive power~\citep{grohe-schweikardt2005}, and gaps of every size are
known between fragments of first-order, temporal, and modal
logics~\citep{etessami-vardi-wilke2002,adler-immerman2003,dawar2007,french2013}.
The closest result to ours concerns finite-variable fragments over a
single binary relation: on finite linear orders, \FOfour\ is
exponentially more succinct than \FOthree, while two- and three-variable
logics are polynomially equivalent~\citep{grohe-schweikardt2005}. Such lower
bounds are proved using the formula-size games of Adler and
Immerman~\citep{adler-immerman2003,hella-vilander2019} or, equivalently, by
tracking a potential function along syntax trees~\citep{grohe-schweikardt2005}.
This framework is indifferent to negation (which merely swaps structure pairs)
and charges a bounded potential increase for each quantifier (corresponding
to composition in \CoR). The core difficulty lies in finding suitable
structures and a potential, neither of which transfers from linear orders to
our setting, for two concrete reasons.

First, on linear orders there is no gap to find: by the lower-bound theorem of
Grohe and Schweikardt~\citep{grohe-schweikardt2005}, an \FOthree\ formula of
size $s$ can only tell in which order its points come and whether two of
them lie at distance exactly $d$ for some $d\le4s^2$, and each such test is
a \CoR\ term of size $O(d)$, namely a $d$-fold composition of
$\mathit{succ}=(<)\cap\overline{(<\,;<)}$. As a formula of size $s$ is a
Boolean combination of the $s^{O(1)}$ available tests, it has an equivalent
\CoR\ term of size $s^{O(1)}$; so \FOthree\ and \CoR\ terms are
polynomially equivalent there.

Second, the mechanism of their gap is fundamentally different. A fourth variable
permits the divide-and-conquer recursion familiar from Savitch's theorem~\citep{savitch1970},
\[
 \varphi_m(x,y) = \exists z\,\forall u\,\bigl((u=x\lor u=y)\to
 \varphi_{m-1}(z,u)\bigr),
\]
which expresses ``$x$ and $y$ are at distance $2^m$'' in size $O(m)$ by
using one copy of $\varphi_{m-1}$ for both halves of the
path~\citep{grohe-schweikardt2005}; with three variables the recursion must
write $\varphi_{m-1}$ twice at every level, and their lower bound shows
that no other \FOthree\ formula does better than $2^{\Omega(m)}$.
Reuse is exactly what a circuit provides: a relational circuit expresses
distance $2^m$ with $m$ compositions, by squaring $S_{i+1}=S_i;S_i$ from
$S_0=\mathit{succ}$, whereas a term needs size $2^{\Omega(m)}$ by their
bound and the linear translation of terms into \FOthree. Our bounds hold for
circuits---where on linear orders the gap actually runs in the opposite
direction---showing that the gap between \FOthree\ and \CoR\ is not simply
one of reuse.

This gap stems from how the two formalisms handle witnesses: in $\Psi_k$, the
quantifier binds $R_i(x,z)\lor R_i(z,y)$ for all $i$ simultaneously,
constraining both edges at once. In contrast, a composition $P;Q$ splits the
edges into $P$ and $Q$. The classical translation \eqref{eq:omega-upper} must
therefore spend one composition on each way of partitioning the $k$ clauses
between the two edges, an overhead that Theorem~\ref{main:omega} proves to be
unavoidable. \citet{nakamura2022} proved the exponential bound for positive
\CoR\ using a formula-size game; however, the argument relies fundamentally on
positivity and does not extend to complement, which exchanges the two sides
and thus requires a symmetric potential. Our preservation theorem supplies
one: for a single finite structure and the family of substructures obtained by
deleting one part each, the potential of a subterm is a set of parts outside
of which every deletion is invisible to it. Boolean operations and converse
contribute nothing to this set, while a composition contributes at most two
parts; because these parts combine by union rather than by addition, a shared
gate is charged only once, directly yielding the bound for circuits. Finally,
$\Psi_k$ and $\Phi_k$, suitably interpreted, detect every deletion.

\paragraph{Composition as matrix multiplication.}
Since composition corresponds to a Boolean matrix product, a \CoR\ term with $m$
compositions evaluates a query using $m$ fast matrix multiplications in time
$O(m\,n^{\omega})$, rather than the $O(n^3)$ time of naive enumeration,
where $\omega$ is the matrix multiplication exponent. Bounded-variable queries
can be evaluated in time polynomial in both the query and the
structure~\citep{vardi1995}. \citet{Williams2014} exploited this speedup for
three-quantifier graph properties, leading to a fine-grained complexity
theory of first-order queries parameterized by variable count and quantifier
structure~\citep{GaoImpagliazzoKolokolovaWilliams2017,BringmannFischerKunnemann2019,GaoImpagliazzo2019}.
The composition counts in Theorems~\ref{main:omega} and~\ref{main:phi}
therefore lower-bound the number of matrix products required by any
relational evaluation of the query, a bound stated directly in
Theorem~\ref{main:matrix}. These results are not lower bounds for general
Boolean circuits with access to individual input entries. In our cost measure,
all Boolean operations (including complement) are free. In contrast,
classical monotone lower bounds for a single Boolean matrix
product~\citep{kerr1970,pratt1975,paterson1975,mehlhorn-galil1976} and
exact bounds for straight-line programs over semirings~\citep{jerrum-snir1982}
count individual gates; see \citet{jukna2012} for a survey.

\paragraph{Relational and matrix query languages.}
Algebraic representations of logical queries have been central to query
processing since Codd~\citep{Codd1970,Codd1972}. Binary relation algebra and its
fragments form the algebraic core of graph and navigational query
languages~\citep{FletcherEtAl2015,HellingsEtAl2021,HellingsEtAl2022}. In these
settings, logic-to-algebra translations can incur exponential
blowups~\citep{BaranyTenCateOtto2012}. The succinctness of relational
representations is therefore a central concern~\citep{FletcherEtAl2015},
because the size of the translated expression directly dictates the cost of
query optimization and evaluation. MATLANG is a matrix query language built
from matrix multiplication, transpose, ones-vectors, diagonalization, and
pointwise functions~\citep{brijder2019tods}; see \citet{geerts2021} for a
survey. For relations annotated with values in a semiring, in the sense of
provenance semirings~\citep{green2007},
\citet{brijder2022} proved that binary-output ARA(3)---the three-attribute
annotated relational algebra---is expressively equivalent to MATLANG,
mirroring the equivalence between \FOthree\ and \CoR. An earlier manuscript
asked whether this exponential arity elimination (and by analogy, the
classical \FOthree\ translation) is unavoidable~\citep[p.~9]{brijder2019}.
Theorem~\ref{main:matrix} resolves this question affirmatively for
commutative idempotent semirings of finite join breadth, while
Section~\ref{semi:positive-section} extends the bound to additional semirings
under a restricted operation set.
